% sections/method.tex -- Section 5 (paper_v2, cut in C4 from the full version's). Every number is
% \val{} from numbers.tex. Figure: figures/f2_ties.tex (F2). B2 (25 November) brought in the
% reviewed Vietnamese version (paper_v2_vn, L10): no opening pointers (F1, the overview, and the
% changes from the base papers go to the supplementary document in B5), the leakage controls and
% the one-standard-error rule as the author reworded them, the runs' agreement in the text, and
% "sampled" for the negatives and LMSPS's paths, so that "draws" means the random sets only. The key
% settings are in Section 6.
\section{Method}
\label{sec:method}

\subsection{Meta-path search for link prediction}
\label{sec:search}

\paragraph{LMSPS\@.} LMSPS \citep{li2024longrange} first computes one feature matrix over the
target nodes for every candidate meta-path: the raw features at the path's end, averaged step by
step along it. A supernet gives each candidate $k$ a weight $\alpha_k$. Each epoch it fuses $M$
paths sampled from a pool with Gumbel-softmax weights computed from $\alpha$. The pool shrinks to
the $V$ paths with the largest $\alpha$. The network's weights learn on training data and $\alpha$
on validation data. Finally, the sets of $M$ paths from the final pool are scored on validation
data, and the best one is the result. We start from the authors'
code\footnote{\url{https://github.com/JHL-HUST/LMSPS}}, which was built for node classification. We
change it for link prediction.

\paragraph{Score and loss.} A product's vector $h_p$ is the weighted sum of its path features (each
through its own linear map), followed by PReLU and dropout. Every user has a free
\val{search-hidden}-dimensional vector $e_u$. A pair is scored by
\begin{equation}
  s(u, p) = e_u^{\top} h_p ,
  \label{eq:score}
\end{equation}
and trained with the BPR loss \citep{rendle2009bpr}, $-\log \sigma\bigl(s(u, p^{+}) - s(u,
p^{-})\bigr)$, over training edges $(u, p^{+})$ and products $p^{-}$ sampled at random from those
$u$ has no training edge to. HAN uses the same score and loss.

\paragraph{Candidates.} The candidates are the paths of 1 to \val{search-hops} hops that start and
end at a product, since HAN needs one product graph per path: \val{search-cands-raw} paths, or
\val{search-cands} once each is merged with its reverse, which gives the same graph once made
symmetric.

\paragraph{Leakage controls.} The features are computed on the training graph, which holds the
training edges themselves. For a training edge, its product's user-path features therefore already
contain the edge's own user, which a held-out edge never has. To avoid this, we learn $\alpha$ on
the validation edges. For the weight step, we compute the user-path features \val{search-folds}
times, each time leaving out a different quarter of the training edges, and read each edge's
features from the version that lacks it.

\paragraph{Schedule and selection.} A set holds $M = 4$ or $M = 6$ paths. The final pool holds
$V = \val{search-pool}$ paths, so that all \val{search-sets} sets of the pool can be scored. A run
lasts \val{search-epochs} epochs, the first \val{search-warmup} with the whole pool. LMSPS scores
sampled sets with fresh Gumbel noise, which on our data moved a set's score more than the gaps
between sets. We score every set once without noise, by its validation BPR against the
\val{split-negatives} fixed negatives of each validation edge. Each run also records its
\emph{tie group}, the sets within $\delta = \val{search-delta}$ of its best.

\subsection{One decision from many runs}
\label{sec:fleet}

One run cannot rank its best sets (Figure~\ref{fig:ties}): its tie group holds a median of
\val{search-tie-med-4} of the \val{search-sets} sets for four paths and \val{search-tie-med-6} for
six, and \val{search-runs} runs of each size produce \val{search-winners-4} and
\val{search-winners-6} different winners. We therefore run the search \val{search-runs} times per
size, with seeds 1 to \val{search-runs} on the same data and split, and take one decision from all
the runs.

% f2_ties.tex -- F2 of the plan: one run cannot rank its top sets. Data: f2_scores.dat and
% f2_picks.dat, written by make_numbers.py from the Fleet v3 records and re-selections.
\begin{figure}[htbp]
  \centering
  \begin{tikzpicture}
    \begin{axis}[
        name=scores, width=\linewidth, height=3.9cm,
        xlabel={Set, best first}, ylabel={Val.\ BPR above the best},
        xmin=0, xmax=211, ymin=0, ymax=0.125, xtick={1,50,100,150,210},
        yticklabel style={/pgf/number format/fixed, /pgf/number format/precision=2},
        scaled y ticks=false, tick label style={font=\scriptsize}, label style={font=\footnotesize},
        title={(a) One run (seed 1): its \val{search-sets} sets}, title style={font=\footnotesize},
        legend style={font=\scriptsize, at={(0.03,0.97)}, anchor=north west, draw=none,
          fill=white, fill opacity=0.85, text opacity=1},
        legend cell align=left, axis line style={black!60}, tick style={black!60}]
      \fill[cSearch!10] (axis cs:0,0) rectangle (axis cs:211,0.010);
      \addplot[cSearch, line width=1pt] table[x=rank, y=gapfour] {figures/f2_scores.dat};
      \addlegendentry{4 paths: \val{search-tie-seed1-4} sets within $\delta$}
      \addplot[cSearch!55, line width=1pt, densely dashed] table[x=rank, y=gapsix]
        {figures/f2_scores.dat};
      \addlegendentry{6 paths: \val{search-tie-seed1-6} sets within $\delta$}
      \draw[black!55, dashed] (axis cs:0,0.010) -- (axis cs:211,0.010)
        node[pos=0.03, above right, font=\scriptsize, text=black!75]
        {tolerance $\delta = \val{search-delta}$};
    \end{axis}
    \begin{axis}[
        at={(scores.below south west)}, anchor=above north west, yshift=-0.35cm,
        width=\linewidth, height=3.3cm,
        ybar, bar width=9pt, ymin=0, ymax=31, ylabel={Distinct sets},
        symbolic x coords={4,6}, xtick=data, xticklabels={4 paths, 6 paths}, enlarge x limits=0.5,
        tick label style={font=\scriptsize}, label style={font=\footnotesize},
        title={(b) Picks over \val{search-runs} runs}, title style={font=\footnotesize},
        legend style={font=\scriptsize, at={(0.5,-0.24)}, anchor=north, draw=none,
          legend columns=1}, legend cell align=left,
        axis line style={black!60}, tick style={black!60}]
      \addplot[fill=cSearch, draw=white] table[x=size, y=own] {figures/f2_picks.dat};
      \addlegendentry{each run's own winner}
      \addplot[fill=cSearch!45, draw=white] table[x=size, y=vote] {figures/f2_picks.dat};
      \addlegendentry{vote across runs}
      \addplot[fill=cOther!45, draw=white, error bars/.cd, y dir=both, y explicit,
          error bar style={black!70, line width=0.6pt}]
        table[x=size, y=nullmed, y error plus=errhi, y error minus=errlo] {figures/f2_picks.dat};
      \addlegendentry{vote on shuffled scores (median, 90\% range)}
    \end{axis}
  \end{tikzpicture}
  \Description{Two panels, one above the other. Top: the validation scores of the 210 path sets of one search run, sorted from best to worst. Many sets lie within the tolerance of the best. Bottom: for four and six paths, bars count the distinct sets that 30 runs pick for each run's own winner, for the vote across runs and for the vote on shuffled scores.}
  \caption{One search run cannot rank its top sets. (a) The validation scores of the
    \val{search-sets} sets one run scores, sorted: many lie within the tolerance $\delta$ of the
    best. (b) Over \val{search-runs} runs, the winners are almost all different. The vote agrees
    more, but only somewhat beyond what it gives on shuffled scores.}
  \label{fig:ties}
\end{figure}


\paragraph{Averaged scores.} The runs score different sets because their final pools differ, but
they rank the sets they share alike (median Spearman correlation \val{search-agree-4} for four paths
and \val{search-agree-6} for six). We fit the validation loss of set $s$ in run $r$ as $a_r + b_s$
by least squares and compare the \val{search-avg-sets-4} and \val{search-avg-sets-6} sets scored in
at least five runs. We then apply a one-standard-error rule: the sets whose $b_s$ lies within one
standard error of the best set's value (\val{search-se-4} for four paths and \val{search-se-6} for
six) form a window, in which the simplest set by total hops comes first, since shorter paths also
make sparser HAN graphs. For four paths this window holds \val{search-window-4} sets: \code{a1} (the
best) and \code{a2}. Both go to the models. For six it holds \val{search-window-6} near-equal sets,
a plateau the search cannot rank. We take the two simplest, as many as for four paths. We also wrote
down the reasons for this decision.

\paragraph{Vote.} A second rule lets each run pick the set in its tie group whose paths the other
runs' winners contain most often. The consensus set \code{c} holds the $M$ paths most common among
the picks. As a vote creates some agreement by construction, we compare it with
\val{search-null-fleets} fleets whose scores are shuffled within each run (Figure~\ref{fig:ties}b).
The runs agree beyond chance, but modestly: \val{search-null-share-4}\% and
\val{search-null-share-6}\% of the shuffled fleets have as few distinct picks.

\paragraph{What goes to the models.} Each size hands over \code{a1}, \code{a2} and \code{c}
(Table~\ref{tab:sets}). The set \code{a1} is the main comparison against the manual sets, fixed in
advance. The tolerance $\delta$ was fixed before the final runs, from the run-to-run noise of
earlier runs. The averaged scores do not depend on it, and a smaller $\delta$ changes neither
\code{a1} nor \code{a2}.

\subsection{Models that use the paths}
\label{sec:models}

\paragraph{HAN\@.} HAN \citep{wang2019heterogeneous} turns each meta-path into a graph over the
target nodes, runs node-level attention (a graph attention layer) on each graph and fuses the
per-path embeddings with semantic-level attention, a weight $\beta$ per path shared by all nodes.
We keep both levels and change three things. \emph{Graphs:} each path's product graph keeps each
product's \val{han-set-topk} strongest neighbours by walk count after every hop, as unpruned paths
with two user hops would link a product to most of the catalogue. The graph is binary and symmetric
and gives every product a self-loop. \emph{Users:} HAN encodes the products only. Every user has a
free vector, and a pair is scored by Eq.~\eqref{eq:score} with BPR. The user vectors hold
\val{han-user-params}M of HAN's \val{han-params-min}--\val{han-params-max}M parameters.
\emph{Path weights:} every path's attention layer starts from the same weights, and $\beta$ is held
equal for the first \val{han-set-beta-warmup} epochs. If each path's attention layer started from its
own random weights, the weight HAN gives a path would depend on the path's position in the list.
Without the warm-up, one path took all the weight within a few epochs.

\paragraph{metapath2vec.} metapath2vec \citep{dong2017metapath2vec} learns node embeddings (without
features) with a skip-gram over random walks that follow a meta-path. From every product we walk
all paths of a set together, each as a repeated template (\code{PUP} gives P--U--P--U\ldots). We
then score the frozen embeddings by Eq.~\eqref{eq:score}. Users are reached only through paths with a
\code{U}, so a set without one leaves the users untrained.

\paragraph{GAT and popularity.} As a baseline without meta-paths, a graph attention network
\citep{velickovic2018graph} trained with BPR runs on the union of all relations, with node types
ignored except for one input projection per type: the homogeneous baseline of the HGB benchmark
\citep{lv2021hgb}. Users and products come out of the same network, so a user is built from the
products the user reviewed. Each training step removes its own positive edges from the graph it
runs on, since a held-out edge is never in the graph. Popularity scores a product by its number of
training users (ties broken at random) and gives every user the same ranking.
